% === Fig 3: participant influence (gen_fig3_influence.py) ===
\begin{figure}[t]
\centering
\includegraphics[width=\textwidth]{figures/fig3_influence.pdf}
\caption{\textbf{The low ICC of unfiltered mean-speed events and the admission contrast were carried by a few participants.}
\textbf{a}, Share of the summed squared session differences of $A$, $\Sigma(S1 - S2)^2$ (R1, floor 5), for the ten largest contributors under
unfiltered mean-speed events in their own cohort ($n$ = 379); the top three (black) account for 87\%.
\textbf{b}, ICC(A,1) after removing the $k$ most influential participants (largest absolute leave-one-out shift in ICC), ranked separately within
each mean-speed construction's own cohort ($n$ before removal: unfiltered 379, centred 377); the unfiltered ICC changed from
.170 to .646 after two removals, and the centred ICC stayed between .727 (no removal) and
.737. Because cohorts and rankings differ, a and b do not isolate an effect of centring.
\textbf{c}, Admission contrast (ICC of admitted minus unfiltered mean-speed events) in the joint cohort ($n$ = 376) after each single-participant omission,
against the omitted participant's share of the summed squared deviations of the session differences from their mean under unfiltered mean-speed events.
Dashed line, full-cohort contrast .319 [.003, .501]; star, the omission that gave .054 [$-$.007, .129]
($n$ = 375); brackets, 95\% participant-cluster bootstrap CIs. \textbf{d}, \textbf{e}, Session 2 against session 1 $A$ in the joint cohort for
unfiltered mean-speed (d, ICC = .374) and admitted unfiltered events (e, ICC = .693) on common axes; star, the participant omitted in c;
dotted line, identity. \textbf{f}, For that participant, the median $l = \ln(\mathrm{PV}_L/\mathrm{PV}_R)$ of return sweeps (RS) and of progressive saccades
in each R1 session under the two constructions; numbers, RS counts. Intervals are descriptive. Participant deletion is a sensitivity analysis here, not a cleaning rule.}
\label{fig:influence}
\end{figure}
